\section{Experimental Details}
\label{app:hyper}
\setcounter{table}{0}
\setcounter{figure}{0}
All experiments use AdamW \cite{loshchilov2019decoupled} with learning rate $10^{-4}$, $\beta_1{=}0.9$, $\beta_2{=}0.999$. No learning rate schedule is applied. Gradients are clipped to a maximum $\ell_2$ norm of 1.0 at every optimizer step. Maximum sequence length is 256 tokens. See Table~\ref{tab:setup} for the complete configuration including ReverseAdaptive defaults.

\subsection{Reproducibility Checklist}
\label{app:repro}
\textbf{Code:} Full source code is publicly released, including all configuration YAML files, training scripts, evaluation scripts, and scripts to reproduce every figure and table from the analysis CSVs.

\textbf{Verification:} A script in the repository checks every number reported in this paper against the released run records.

\textbf{Seeds:} The IID seeds are listed in Table~\ref{tab:setup}. The $\alpha{=}0.1$ partition-sensitivity runs additionally use seeds 789 and 1000, which appear only in Table~\ref{tab:b2}. Per-seed results are in Appendix~\ref{app:perseed}.

\textbf{Hardware:} 34 runs were executed on Apple silicon (MPS), a single Apple M4 Pro workstation with 48\,GB unified memory, approximately 285 hours of compute, and re-executed on commercially rented NVIDIA hardware, an RTX 4090 for TinyLlama-1.1B and an A100 40\,GB for LLaMA-3.2-3B; the comparison is in Appendix~\ref{app:crossbackend}. The Dolly-15k replication and the FFA-LoRA and FedIT baselines were produced on the rented hardware only.

\textbf{Data:} Alpaca \cite{taori2023alpaca} via \texttt{tatsu-lab/alpaca} on HuggingFace, 3000-sample subset, and Dolly-15k \cite{dolly2023} via \texttt{databricks/databricks-dolly-15k}, matched subset. Adapter checkpoints are not released; reproduction scripts are sufficient.

\subsection{Disabling the Switch}
\label{app:noswitch}
The no-switch sanity baseline deserves specific comment. Running ReverseAdaptive with switching disabled reproduces the matched FLoRA \cite{wang2024flora} run to 10 decimal places at every round, establishing that the wrapper is a zero-cost modification when unused: it rules out numerical artifacts from the wrapper's presence, not merely from its activation. The comparison is within a single backend and a single execution environment. Agreement across hardware is weaker than bit-identity and is reported in Appendix~\ref{app:crossbackend}.

Disabling the switch requires care. Setting $\tau$ to a small positive value does not achieve it: any round in which the training loss fails to improve satisfies $\rho_r < \tau$, so the criterion fires. On the FLoRA trajectory underlying the threshold ablation the loss rises once, at round 11, so $\tau{=}0$ still triggers a switch at that round, and suppressing the switch entirely requires $\tau$ below the most negative round-over-round relative change observed, approximately ${-}0.0009$ here. The sanity baseline accordingly uses $\tau{=}{-}1.0$ together with a warmup exceeding the round budget. A practitioner wanting the wrapper present but inert should disable it by one of those means rather than by choosing a conservative threshold.

\needspace{3\baselineskip}
\section{Per-Seed Full Results}
\label{app:perseed}
\setcounter{table}{0}
\setcounter{figure}{0}
Tables in this appendix are labeled with their corpus, which follows the corresponding table in the main text. The TinyLlama-1.1B method comparisons use the CUDA corpus, since the FFA-LoRA, FedIT, and Dolly-15k runs exist only there; the LLaMA-3.2-3B results use MPS, matching Table~\ref{tab:llama_iid}; and the partition-sensitivity results retain the larger MPS seed set, where five seeds at $\alpha{=}0.1$ provide better evidence than three.

\begin{table*}[tbp]
\centering
\footnotesize
\scriptsize
\renewcommand{\arraystretch}{0.78}
\setlength{\tabcolsep}{6pt}
\caption{Per-seed TinyLlama-1.1B IID results, CUDA corpus, backing Table~\ref{tab:frontier}. Communication is deterministic across seeds. $\dagger$ For Two-Phase the Switch/boundary column shows the predetermined phase boundary $K$, not an adaptive switch event; for FFA-LoRA, $A$ is frozen from initialization (shown as ``init''); for ReverseAdaptive it shows the round at which the plateau signal triggered}
\label{tab:b1}
\begin{tabular}{@{}llrrr@{}}
\toprule
Method & Seed & Final loss & Total MB & Switch / boundary \\
\midrule
FLoRA & 42 & 1.261408 & 2578.13 & -- \\
FLoRA & 123 & 1.260479 & 2578.13 & -- \\
FLoRA & 456 & 1.260500 & 2578.13 & -- \\
FedIT & 42 & 1.260767 & 2578.13 & -- \\
FedIT & 123 & 1.258420 & 2578.13 & -- \\
FedIT & 456 & 1.261503 & 2578.13 & -- \\
Two-Phase $K{=}8$ & 42 & 1.271276 & 1863.13 & $K{=}8$ \\
Two-Phase $K{=}8$ & 123 & 1.270130 & 1863.13 & $K{=}8$ \\
Two-Phase $K{=}8$ & 456 & 1.270237 & 1863.13 & $K{=}8$ \\
ReverseAdaptive ($\tau{=}0.01$) & 42 & 1.275514 & 1533.13 & 6 \\
ReverseAdaptive ($\tau{=}0.01$) & 123 & 1.274610 & 1533.13 & 6 \\
ReverseAdaptive ($\tau{=}0.01$) & 456 & 1.274481 & 1533.13 & 6 \\
FFA-LoRA & 42 & 1.303975 & 983.13 & init \\
FFA-LoRA & 123 & 1.301738 & 983.13 & init \\
FFA-LoRA & 456 & 1.303628 & 983.13 & init \\
\bottomrule
\end{tabular}
\end{table*}

At $\alpha{=}0.5$, mean final loss is statistically indistinguishable from IID for every method in the released analysis CSVs, with differences of 0.001 to 0.007 against per-seed standard deviations near 0.02 over three seeds. Heterogeneity manifests instead as sensitivity to the partition draw: seed-to-seed standard deviation rises by roughly an order of magnitude, from $0.5$ to $1.6 \times 10^{-3}$ under IID in Table~\ref{tab:b1} to $18$ to $25 \times 10^{-3}$ at $\alpha{=}0.5$ across the methods in the released CSVs, which is also why the per-seed spread at $\alpha{=}0.1$ in Table~\ref{tab:b2} is wide.

\begin{table*}[tbp]
\centering
\footnotesize
\setlength{\tabcolsep}{6pt}
\caption{Non-IID results, MPS corpus, five seeds at $\alpha{=}0.1$. This table documents partition sensitivity rather than comparing methods, so the larger MPS seed set is retained; Table~\ref{tab:b1} follows the CUDA corpus because it backs a method comparison. Seeds 42 and 1000 at $\alpha{=}0.1$ had 9 and 8 active clients due to empty Dirichlet partitions, which scales communication to $9/10$ and $8/10$ of the full value (1379.81 and 1226.50 MB against 1533.13). Cross-backend verification (Appendix~\ref{app:crossbackend}) covers seeds 42, 123, and 456; seeds 789 and 1000 were executed on MPS only}
\label{tab:b2}
\begin{tabular}{@{}llrrrr@{}}
\toprule
Setting & Seed & Final loss & Total MB & Switch round & Active clients \\
\midrule
ReverseAdaptive, $\alpha{=}0.5$ & 42 & 1.285933 & 1533.13 & 6 & 10 \\
ReverseAdaptive, $\alpha{=}0.5$ & 123 & 1.282625 & 1533.13 & 6 & 10 \\
ReverseAdaptive, $\alpha{=}0.5$ & 456 & 1.252133 & 1533.13 & 6 & 10 \\
ReverseAdaptive, $\alpha{=}0.1$ & 42 & 1.268353 & 1379.81 & 6 & 9 \\
ReverseAdaptive, $\alpha{=}0.1$ & 123 & 1.288985 & 1533.13 & 6 & 10 \\
ReverseAdaptive, $\alpha{=}0.1$ & 456 & 1.348083 & 1533.13 & 6 & 10 \\
ReverseAdaptive, $\alpha{=}0.1$ & 789 & 1.259200 & 1533.13 & 6 & 10 \\
ReverseAdaptive, $\alpha{=}0.1$ & 1000 & 1.273046 & 1226.50 & 6 & 8 \\
\bottomrule
\end{tabular}
\end{table*}

\begin{table*}[tbp]
\centering
\footnotesize
\setlength{\tabcolsep}{6pt}
\caption{Per-seed LLaMA-3.2-3B IID results, MPS corpus, backing Table~\ref{tab:llama_iid}. ReverseAdaptive uses $\tau{=}0.01$. Communication is deterministic given the switch round; the three ReverseAdaptive seeds occupy three adjacent lattice points}
\label{tab:b3}
\begin{tabular}{@{}llrrr@{}}
\toprule
Method & Seed & Final loss & Total MB & Switch \\
\midrule
FLoRA & 42 & 1.371683 & 2625.0 & -- \\
FLoRA & 123 & 1.373402 & 2625.0 & -- \\
FLoRA & 456 & 1.371998 & 2625.0 & -- \\
Two-Phase $K{=}8$ & 42 & 1.393776 & 1942.5 & $K{=}8$ \\
Two-Phase $K{=}8$ & 123 & 1.394877 & 1942.5 & $K{=}8$ \\
Two-Phase $K{=}8$ & 456 & 1.392772 & 1942.5 & $K{=}8$ \\
ReverseAdaptive & 42 & 1.393831 & 1837.5 & 8 \\
ReverseAdaptive & 123 & 1.390265 & 1942.5 & 9 \\
ReverseAdaptive & 456 & 1.397366 & 1732.5 & 7 \\
\bottomrule
\end{tabular}
\end{table*}

\begin{table*}[tbp]
\centering
\footnotesize
\setlength{\tabcolsep}{6pt}
\caption{Per-seed TinyLlama-1.1B Dolly-15k IID results, CUDA corpus, backing Table~\ref{tab:dolly}. Held-out $\Delta$ is tuned minus base on the Dolly-15k held-out slice; more negative is better}
\label{tab:b4}
\begin{tabular}{@{}llrrr@{}}
\toprule
Method & Seed & Final loss & Held-out $\Delta$ & Total MB \\
\midrule
FLoRA & 42 & 1.656762 & $-0.531848$ & 2578.13 \\
FLoRA & 123 & 1.652934 & $-0.533376$ & 2578.13 \\
FLoRA & 456 & 1.653426 & $-0.533664$ & 2578.13 \\
Two-Phase $K{=}8$ & 42 & 1.666616 & $-0.527571$ & 1863.13 \\
Two-Phase $K{=}8$ & 123 & 1.662685 & $-0.529113$ & 1863.13 \\
Two-Phase $K{=}8$ & 456 & 1.663475 & $-0.529410$ & 1863.13 \\
ReverseAdaptive & 42 & 1.671039 & $-0.526568$ & 1533.13 \\
ReverseAdaptive & 123 & 1.666995 & $-0.527418$ & 1533.13 \\
ReverseAdaptive & 456 & 1.667699 & $-0.526547$ & 1533.13 \\
\bottomrule
\end{tabular}
\end{table*}

\section{Statistical Tests}
\label{app:stats}
\setcounter{table}{0}
\setcounter{figure}{0}
Paired $t$-tests over three seeds on final-round training loss and on held-out instruction-following loss. TinyLlama-1.1B results use the CUDA corpus, matching Table~\ref{tab:frontier}; LLaMA-3.2-3B results use the MPS corpus.

With three seeds these tests have very low power, so effect sizes should be read alongside the $p$-values. No multiplicity correction is applied in the main results tables. Applying a Bonferroni correction across the six TinyLlama-1.1B Alpaca IID comparisons gives a threshold of $0.0083$, which the ReverseAdaptive vs.\ Two-Phase $K{=}8$ comparison on held-out loss ($p = 0.026$) does not meet. That comparison is secondary under the frontier framing of Section~\ref{sec:discussion}; the load-bearing comparison is ReverseAdaptive vs.\ FFA-LoRA, which holds at $p < 10^{-3}$ on both metrics with a difference roughly twenty times the largest per-method standard deviation.

The LLaMA-3.2-3B comparison between ReverseAdaptive and Two-Phase $K{=}8$ has a 95\% confidence interval of $[{-}0.0114, 0.0114]$, which contains the $0.0043$ difference observed between the same methods at TinyLlama-1.1B. The test therefore cannot distinguish equivalence from a difference of the magnitude seen at the smaller scale.

\begin{table*}[tbp]
\centering
\footnotesize
\setlength{\tabcolsep}{5pt}
\renewcommand{\arraystretch}{0.88}
\caption{Paired $t$-tests, three seeds. TinyLlama-1.1B rows use the CUDA corpus; LLaMA-3.2-3B rows use MPS. Positive $\Delta$ indicates the first method has higher loss. Held-out instruction-following loss was not computed for the LLaMA-3.2-3B corpus; quality at that scale is reported as final training loss and downstream zero-shot accuracy}
\label{tab:c1}
\begin{tabular}{@{}lrrrr@{}}
\toprule
Comparison & $\Delta$ final loss & $p$ & $\Delta$ held-out & $p$ \\
\midrule
\multicolumn{5}{l}{\emph{TinyLlama-1.1B, Alpaca, IID}} \\
Two-Phase $K{=}8$ vs.\ FLoRA & $+0.0098$ & $4.2 \times 10^{-5}$ & $+0.0034$ & $7.9 \times 10^{-5}$ \\
ReverseAdaptive vs.\ FLoRA & $+0.0141$ & $1.1 \times 10^{-5}$ & $+0.0063$ & $6.2 \times 10^{-3}$ \\
ReverseAdaptive vs.\ Two-Phase $K{=}8$ & $+0.0043$ & $3.4 \times 10^{-4}$ & $+0.0029$ & $2.6 \times 10^{-2}$ \\
FFA-LoRA vs.\ ReverseAdaptive & $+0.0282$ & $4.4 \times 10^{-4}$ & $+0.0182$ & $3.8 \times 10^{-4}$ \\
FFA-LoRA vs.\ FLoRA & $+0.0423$ & $1.7 \times 10^{-4}$ & $+0.0245$ & $1.1 \times 10^{-4}$ \\
FedIT vs.\ FLoRA & ${-}0.0006$ & $0.588$ & $+0.0002$ & $0.425$ \\
\midrule
\multicolumn{5}{l}{\emph{TinyLlama-1.1B, Dolly-15k, IID}} \\
Two-Phase $K{=}8$ vs.\ FLoRA & $+0.0099$ & $7.8 \times 10^{-5}$ & $+0.0043$ & $2.4 \times 10^{-6}$ \\
ReverseAdaptive vs.\ FLoRA & $+0.0142$ & $2.5 \times 10^{-5}$ & $+0.0061$ & $7.6 \times 10^{-3}$ \\
\midrule
\multicolumn{5}{l}{\emph{LLaMA-3.2-3B, Alpaca, IID}} \\
Two-Phase $K{=}8$ vs.\ FLoRA & $+0.0214$ & $3.2 \times 10^{-4}$ & -- & -- \\
ReverseAdaptive vs.\ FLoRA & $+0.0215$ & $1.3 \times 10^{-2}$ & -- & -- \\
ReverseAdaptive vs.\ Two-Phase $K{=}8$ & $+0.000012$ & $0.997$ & -- & -- \\
\bottomrule
\end{tabular}
\end{table*}

\section{Cross-Backend Verification}
\label{app:crossbackend}
\setcounter{table}{0}
\setcounter{figure}{0}

The primary corpus was produced on the MPS backend and subsequently re-executed on CUDA. Table~\ref{tab:d1} summarizes the comparison; Table~\ref{tab:d2} details every disagreement.

Tolerances are setting-aware. IID runs require per-round training loss to agree within $0.01$ absolute. Non-IID runs allow $0.025$, since heterogeneous partitions amplify per-round variation, but additionally require the switch round to match exactly and permit no more than two consecutive rounds exceeding $0.01$. Communication must agree within $0.01$ in all settings, which is effectively exact since totals are protocol-deterministic.

Of 34 run pairs, 31 have a valid MPS reference. The three Two-Phase $K{=}8$ runs at $\alpha{=}0.5$ do not: their MPS executions predate the bidirectional B-only implementation and never switched, recording 2578.13 MB, the full FLoRA volume. Those are absent references rather than disagreements, and CUDA is the source of truth for that configuration. Of the 31 comparable pairs, 27 agree within tolerance and 4 do not.

\begin{table}[tbp]
\centering
\footnotesize
\renewcommand{\arraystretch}{0.88}
\setlength{\tabcolsep}{6pt}
\caption{Cross-backend verification summary, MPS against CUDA, 34 run pairs}
\label{tab:d1}
\begin{tabular}{@{}lr@{}}
\toprule
Category & Runs \\
\midrule
Agreed within setting-aware tolerance & 27 \\
Disagreed & 4 \\
\midrule
Comparable pairs & 31 \\
No valid MPS reference (CUDA reported) & 3 \\
\midrule
Total & 34 \\
\bottomrule
\end{tabular}
\end{table}

\begin{table}[tbp]
\centering
\footnotesize
\scriptsize
\renewcommand{\arraystretch}{0.85}
\setlength{\tabcolsep}{3pt}
\caption{Cross-backend disagreements. Every communication difference is one 110\,MB lattice step. Run 16 failed only the consecutive-spike rule}
\label{tab:d2}
\begin{tabular}{@{}llrrr@{}}
\toprule
Run & Config. & MPS & CUDA & Note \\
\midrule
11 & $\tau{=}0.001$, IID & 2083.13 & 1973.13 & $s{=}11$/$10$ \\
12 & $\tau{=}0.002$, IID & 1863.13 & 1753.13 & $s{=}9$/$8$ \\
16 & $\alpha{=}0.1$, 123 & 1533.13 & 1533.13 & spike rule \\
17 & $\alpha{=}0.1$, 456 & 1533.13 & 1643.13 & $s{=}6$/$7$ \\
\bottomrule
\end{tabular}
\end{table}

Three of the four disagreements are the same phenomenon: every observed cross-backend communication difference is exactly one 110\,MB switch-round step (Section~\ref{sec:measuredbytes}), so Table~\ref{tab:d2} lists adjacent lattice points. The backends agree on the protocol and disagree only on which round the plateau fires, and only near a boundary---two tight-$\tau$ IID runs ($\tau{=}0.001$, $\tau{=}0.002$) and one $\alpha{=}0.1$ run. Run 16 has no communication or switch-round difference: only the consecutive-spike rule failed.

Run 17 was re-executed on CUDA with partition seed 456 and training seed 999 and again switched at round 7, so the one-round gap is a stable backend/partition property rather than run noise. At $\alpha{=}0.1$, exact switch-round agreement is therefore not guaranteed (one of three CUDA counterparts fires one round later); CUDA is reported as primary. The switch round matched exactly for every IID run at $\tau{=}0.01$ and for all three LLaMA-3.2-3B ReverseAdaptive runs, supporting the cross-scale claim in Section~\ref{sec:choosing}.

\textbf{Byte accounting.} Transport-layer counts use \texttt{numel}$\times$\texttt{element\_size} on tensors actually sent; the FFA-LoRA \texttt{get\_communication\_cost} helper is unused.

\textbf{Code revisions.} Runs span multiple revisions (including uncommitted working-tree changes), so numbers come from released records and the replication here. Byte totals are revision-invariant: FLoRA, Two-Phase $K{=}8$ and ReverseAdaptive each matched under two CUDA revisions at 2578.13, 1863.13 and 1533.13\,MB. Loss carries revision uncertainty; the FLoRA and FedIT rows of Table~\ref{tab:frontier} differ by 0.0006 across that boundary, against a ReverseAdaptive-to-FFA-LoRA gap of 0.0282.
